% 本文件正文由作者撰写。
    本研究使用了三个数据集,其中2和3是MEG,1是EEG,1和2是中文,3是西班牙语.三者的任务范式都是被动聆听,1和3聆听的都是
连续故事<<小王子>>,2聆听的是句子.
    三个数据集都只做了简单的预处理:0.1-40hz的带通滤波,重采样到50hz.
    本研究使用1s的统一词级窗口,从词的oneset开始向后截取1s的神经活动,由于1数据集聆听的是实际朗读的音频,我们人工对
实际朗读的音频进行了审查,并使用强制对齐工具获得了字符级和词级时间戳.
    在数据集的划分上,对于1和3按照章节来划分,2按照领域来划分.过往的研究中,有采用句子内容哈希的形式进行划分.对于我们
的研究,之所以采用章节划分的原因是因为,我们发现1数据集总是会学到缓慢的试次漂移现象.
    脑电解码的模型我们使用了膨胀卷积提取字符窗口的特征,并将提取后的词级特征输入句子级Transformer模型中以利用连续语言
神经活动.对于1数据集,由于聆听任务中受试者聆听的是该数据集朗读任务范式的真实音频,而朗读任务中刺激材料是按照0.25s的视觉
高亮时钟进行逐行显示的,部分句子内容会被截断,所以我们将
